%
%
%

\begin{frame}[t,allowframebreaks]{Applying a linear model to the XOR regression task -} 

    We can choose a linear form for the model $f(\mathbf{x};\mathbf{\theta})$,
    where the parameter set $\mathbf{\theta}$ consists 
    of a set of weights $\mathbf{w}$ = $(w_1,w_2)$ and a bias $b$.        

    \begin{columns}[t]
        \begin{column}{0.50\textwidth}
            \vspace{-0.6cm}
            \begin{center}
                \begin{tikzpicture}[scale=1]
                  %\draw[help lines] (0,0) grid (6,3.7);
                  \node[ann_processing_node] (o)  at (3.0, 1.5) {$\sum$};
                  \node[ann_input_node]      (x1) at (0.0, 3.0) {$x_1$};
                  \node[ann_input_node]      (x2) at (0.0, 1.5) {$x_2$};
                  \node[ann_bias_node]       (b)  at (1.0, 0.0) {$+1$};
              
                  \drawgraphlinebigarrow (x1.east) 
                  to node[above, midway] 
                  {\small $w_1$}(o.north west) ;
              
                  \drawgraphlinebigarrow (x2.east) 
                  to node[above, midway] 
                  {\small $w_2$}(o.west) ;
              
                  \drawgraphlinebigarrow (b.east) 
                  to node[above, midway] 
                  {\small $b$}(o.south west) ;
              
                  \drawgraphlinebigarrow (o.east) 
                  to node[above,midway] 
                  {\small \color{black} $f(\mathbf{x};\mathbf{w},b)$} (6.0,1.5);              
                \end{tikzpicture}
            \end{center}        
        \end{column}
        \begin{column}{0.50\textwidth}
            \begin{equation}
                f(\mathbf{x};\mathbf{w},b) = \mathbf{x}^{T} \mathbf{w} + b
                \label{eq:learn_xor_linear_model_1}
            \end{equation}

            As discussed previously, we can rewrite Eq.~\label{eq:learn_xor_linear_model_1}
            in a more compact form by adding an additional input $x_0$ that always takes the value of +1 
            and considering the bias $b$ to be the weight $w_0$
            \begin{equation}
                f(\mathbf{x};\mathbf{w}) = \mathbf{x}^{T} \mathbf{w}
                \label{eq:learn_xor_linear_model_2}
            \end{equation}
        \end{column}
    \end{columns}
      
    \framebreak

    %
    %

    The \index{loss function}\gls{loss function} 
    of Eq.~\ref{eq:learn_xor_loss_function_1} 
    could be minimised using the 
    \index{gradient descent}\gls{gradient descent} 
    method we saw in Part \prevpart, 
    and which will be further discussed in Part \nextpart.\\
    \vspace{0.1cm}
    Owing to the simplicity of the given 
    \index{regression}\index{linear regression}\gls{regression} problem, 
    it is possible to obtain a solution in closed form 
    (see derivation of Eq.~\ref{eq:linear_regression_closed_form_solution_matrix}).\\
    \vspace{-0.2cm}
    \begin{equation*}
        \vect{w} = (\vect{X}^T \vect{X})^{-1} \vect{X}^T \vect{y}
    \end{equation*}\\
    \vspace{0.1cm}
    The training set $\mathbb{D}$ contains 4 examples 
    $(\mathbf{x}^{(i)},y^{(i)})$, with $i \in [1,4]$, given in the table below.\\
    \vspace{0.1cm}
    \begin{columns}[t]
        \begin{column}{0.44\textwidth}
            \vspace{-0.6cm}
            \begin{center}
                \begin{tabular}{ c | c c c | c }
                 $i$ & $x_0$ & $x_1$ & $x_2$ & $y = x_1 \oplus x_2$ \\ 
                 \hline
                 1 & 1 & 0 & 0 & 0 \\  
                 2 & 1 & 0 & 1 & 1 \\   
                 3 & 1 & 1 & 0 & 1 \\  
                 4 & 1 & 1 & 1 & 0 \\   
                \end{tabular}
            \end{center}
        \end{column}
        \begin{column}{0.56\textwidth}
            The inputs $\mathbf{x}^{(i)}$ and target values
            $y^{(i)}$ spanning $\mathbb{D}$ can be packed into matrices
            as shown in Eq.~\ref{eq:linear_regression_closed_form_solution_matrix_inputs}.
            \begin{equation}
                \vect{X} = 
                \begin{pmatrix}
                    1 & 0 & 0 \\
                    1 & 0 & 1 \\
                    1 & 1 & 0 \\
                    1 & 1 & 1 \\
                \end{pmatrix} 
                \; \textrm{and} \;
                \vect{y} = 
                \begin{pmatrix}
                    0 \\
                    1 \\
                    1 \\
                    0 \\
                \end{pmatrix} 
            \end{equation}        
        \end{column}
    \end{columns}

    \framebreak

    %
    %

    The matrix product $\vect{X}^T \vect{X}$ is:\\
    \vspace{-0.3cm}
    \begin{equation}
        \vect{X}^T \vect{X} = 
        \begin{pmatrix}
            1 & 1 & 1 & 1 \\
            0 & 0 & 1 & 1 \\
            0 & 1 & 0 & 1 \\
        \end{pmatrix} 
        \begin{pmatrix}
            1 & 0 & 0 \\
            1 & 0 & 1 \\
            1 & 1 & 0 \\
            1 & 1 & 1 \\
        \end{pmatrix} =
        \begin{pmatrix}
            4 & 2 & 3 \\
            2 & 2 & 2 \\
            2 & 1 & 2 \\
        \end{pmatrix} 
    \end{equation}        

    while its inverse, as it can be easily verified, is:\\
    \vspace{-0.3cm}
    \begin{equation}
        (\vect{X}^T \vect{X})^{-1} = 
        \begin{pmatrix}
            0 & -0.5 & -1 \\
            0 &  1   & -1 \\
           -1 &  0   &  2 \\
        \end{pmatrix} 
    \end{equation}        
    \vspace{-0.5cm}
    Substituting the above matrices 
    in Eq.~\ref{eq:linear_regression_closed_form_solution_matrix}, we find:\\
    \begin{equation}
        \vect{w} 
        = 
        \begin{pmatrix}
            b   \\
            w_1 \\
            w_2 \\
        \end{pmatrix} 
        =
        \begin{pmatrix}
            0 & -0.5 & -1 \\
            0 &  1   & -1 \\
           -1 &  0   &  2 \\
        \end{pmatrix} 
        \begin{pmatrix}
            1 & 1 & 1 & 1 \\
            0 & 0 & 1 & 1 \\
            0 & 1 & 0 & 1 \\
        \end{pmatrix} 
        \begin{pmatrix}
            0 \\
            1 \\
            1 \\
            0 \\
        \end{pmatrix} 
        =
        \begin{pmatrix}
            0.5 \\
            0 \\
            0 \\
        \end{pmatrix} 
        \label{eq:learn_xor_linear_model_weights}
    \end{equation}        

    \framebreak

    %
    %

    The estimated model parameters $\mathbf{w}^T=(0.5, 0, 0)$,
    substituted in Eq.~\ref{eq:learn_xor_linear_model_2}, 
    yield the following predictions:\\
    \vspace{0.4cm}
    \begin{equation}
        \hat{y}(\mathbf{x}) = 
        \begin{pmatrix}
            1 & x_1 & x_2 \\
        \end{pmatrix} 
        \begin{pmatrix}
            b   \\
            w_1 \\
            w_2 \\
        \end{pmatrix} =
        \begin{pmatrix}
            1 & x_1 & x_2 \\
        \end{pmatrix} 
        \begin{pmatrix}
            0.5 \\
            0   \\
            0   \\
        \end{pmatrix} = 0.5
        \label{eq:learn_xor_linear_model_predictions}
    \end{equation}

    Therefore, irrespective of the input $\mathbf{x}$, the {\bf model always returns 0.5}, 
    halfway between the two possible target values of 0 and 1.\\
    \vspace{0.2cm}
    This is a rather {\bf spectacular failure} for our trained linear model.

\end{frame}

%
%
%

\begin{frame}[t,allowframebreaks]{Linear models cannot learn the XOR function -} 

    \begin{columns}[t]
        \begin{column}{0.50\textwidth}
            \begin{center}
                \includegraphics[width=0.90\textwidth]
                    {./images/people/minsky_1.png}\\
                {\scriptsize 
                Marvin Misky. 
                \color{col:attribution} 
                Photograph by Robert Kaiser (Associated Press).\\}        
            \end{center}
        \end{column}
        \begin{column}{0.50\textwidth}
        \end{column}
    \end{columns}


\end{frame}
